
This work tests whether the curvature of the ERM loss surface --- measured as per-sample Hessian trace of the last fully-connected layer --- reveals minority group membership without group annotations. It does: 4/5 random seeds achieve $AUROC\geq 0.85$ at the epoch of maximum trace ratio, with the epoch-0 control confirming that the signal is created by ERM training on Waterbirds, not inherited from ImageNet pretraining (epoch-0 AUROC $\in$ [0.538, 0.609] in all five seeds). The Hutchinson estimator at K=50 is stable across seeds (CV<3\%), establishing this as the minimum reliable setting for last-fc of ResNet-50.

The mechanism investigation produces an equally clear result, though in the negative direction. The pre-registered mechanism gate --- minority training-set confidence remaining in [0.3, 0.7] at t\textit{ --- fails in all five seeds (p\_minority(t})=0.9678--0.9999). A transient early-epoch confidence differential does exist at epoch 1 (seed 1: p\_min=0.827 vs p\_maj=0.972), consistent with Phase I theory \citep{labonte2026phase}, but it disappears before t\textit{ in all seeds. The confidence channel p\_i($1-p_{i})$ is inactive at t} for both groups.

The remaining candidate for the trace asymmetry at t\textit{ is the feature-norm channel: systematic elevation of $\|x_{i}\|^2$ for minority samples, consistent with LaBonte et al. [2024]'s group-level spectral imbalance result. Directly verifying this --- measuring penultimate-layer norms at t} and correlating them with trace ranks --- is the primary next mechanistic step. The practical path forward requires testing whether the trace proxy, used to guide annotation-free last-layer retraining (DFR-style), improves worst-group accuracy on Waterbirds.

The existence result establishes the signal quality that makes the downstream application worth testing; the mechanism result identifies what signal the proxy is likely capturing. A second-order perspective on spurious feature reliance --- one that asks how ERM differentially shapes the loss landscape curvature around training samples, not merely how it classifies them --- opens a class of annotation-free minority detectors orthogonal to first-order signals.

\subsubsection{Broader Impact}

This work advances annotation-free methods for identifying minority groups in machine learning training data. The ability to detect spuriously correlated minority samples without group labels could reduce the annotation burden for fairness-motivated model corrections. However, the same signal could in principle be used to identify and selectively filter minority samples in adversarial settings. Application of this method should be in contexts where minority detection supports, rather than undermines, equitable model performance.
